% Gerado por src/gerar_tabelas.py a partir de data/results/ -- não edite à mão.
% Requer \usepackage{booktabs}.
% A matriz de confusão por janela da validação cruzada do modelo base (treino original) não foi gravada em
% nenhum arquivo de resultado (train_classifier.py só gerou data/results/confusion_matrix.png).

\begin{table}[htbp]
\centering
\small
\caption{Comparação registrada entre o modelo base treinado novamente com sementes fixas e o modelo com rótulos por intervalo}
\label{tab:comparacao}
\begin{tabular}{lrrrr}
\toprule
 & \multicolumn{2}{c}{Métrica 1} & \multicolumn{2}{c}{Métrica 2} \\
\cmidrule(lr){2-3}\cmidrule(lr){4-5}
 & Base & Intervalos & Base & Intervalos \\
\midrule
Agregado (56 vídeos) & 0,562 & 0,322 & 0,020 & 0,014 \\
\midrule
Partição 1 & 0,501 & 0,177 & 0,011 & 0,003 \\
Partição 2 & 0,455 & 0,328 & 0,018 & 0,023 \\
Partição 3 & 0,505 & 0,471 & 0,035 & 0,012 \\
Partição 4 & 0,596 & 0,398 & 0,016 & 0,018 \\
Partição 5 & 0,582 & 0,307 & 0,025 & 0,014 \\
Partições com Intervalos maior & \multicolumn{2}{c}{0 de 5} & \multicolumn{2}{c}{2 de 5} \\
\midrule
\multicolumn{4}{l}{Melhora nas duas métricas agregadas} & não \\
\multicolumn{4}{l}{Melhora nas duas métricas em partições suficientes, pelo critério} & não \\
\multicolumn{5}{l}{Veredito: sem diferença distinguível} \\
\bottomrule
\end{tabular}
\legend{Linha ``Agregado'': valores agregados sobre os 56 vídeos; demais linhas: valor de cada partição. Métrica 1: macro F1 por quadro, restrito aos quadros anotados. Métrica 2: F1 por evento, IoU 0,3, regra principal. Base: modelo base treinado novamente com sementes fixas, com rótulo por vídeo; Intervalos: modelo com rótulos por intervalo anotado e classe de fundo. Critério registrado antes da avaliação (criterio\_comparacao.md): os rótulos por intervalo só seriam considerados melhores se as duas condições acima fossem atendidas. Fonte: Elaborado pelo autor.}
\end{table}
